\documentclass{article}
\usepackage{neurips_2024}
\usepackage[utf8]{inputenc}
\usepackage[T1]{fontenc}
\usepackage{hyperref}
\usepackage{url}
\usepackage{booktabs}
\usepackage{multirow}
\usepackage{amsfonts}
\usepackage{amsmath}
\usepackage{amssymb}
\usepackage{nicefrac}
\usepackage{microtype}
\usepackage{xcolor}
\usepackage{listings}

\newcommand{\sio}{\mathbb{S}}
\newcommand{\ZD}{\mathrm{ZD}}
\newcommand{\ExactlyPrivate}[1]{\texttt{ExactlyPrivate}\langle #1 \rangle}
\newcommand{\Editable}[1]{\texttt{Editable}\langle #1 \rangle}
\newcommand{\CapabilityGated}[1]{\texttt{CapabilityGated}\langle #1 \rangle}
\newcommand{\Forgettable}[1]{\texttt{Forgettable}\langle #1 \rangle}
\newcommand{\withZD}{\texttt{with ZD}}

\title{Algebraically Exact Surgical Interventions on Trained Models\\via Sedenion Zero-Divisor Algebra}
\author{%
  Anonymous Author(s)\\
  Paper under double-blind review%
}

\begin{document}
\maketitle

\begin{abstract}
Three open problems in applied machine learning --- data-subject removal
(GDPR/RTBF), ripple-free model editing, and AI-safety capability containment
--- currently admit only \emph{approximate} solutions whose guarantees depend
on epsilon-bounds, rank-1 projections, or mechanistic-interpretability
heuristics.  We observe that all three are instances of a single algebraic
schema: \textsc{Erase}(subspace) $\land$ \textsc{Preserve}(complement).  When
the subspace is chosen as the 4-dimensional right-annihilator of a primitive
sedenion zero-divisor operator $A$, the operation becomes
\emph{algebraically exact}.  Building on the 168-class Bridge Theorem (Paper~B;
\texttt{SounioZeroDivisorBridge.lean}), we introduce three compiler-level
wrapper types --- $\ExactlyPrivate{T}$, $\Editable{T}$, $\CapabilityGated{T}$
--- each enforcing a \withZD~effect at use-sites and backed by a Lean~4
theorem for the underlying algebraic guarantee.  Three Sounio benchmarks
demonstrate the operational consequences: exact ($0.000000$ L2) unlearning
against classical methods' residuals ($0.008$ gradient-ascent, $0.78$
differential-privacy shrinkage); $0$ ripple on ZD-kernel-bound edits against
$0.51$ leakage for rank-1 ROME-style updates; and danger-capability accuracy
falling from $100\%$ to chance $(50.5\%)$ under the ZD erase while safe
capability is preserved at $100\%$.  The contribution is not a new neural
architecture: it is a \emph{type-level} guarantee grounded in structural
algebra that previous work could not articulate because the division-algebra
substrate (reals, complexes, quaternions, octonions) does not carry
zero-divisors.
\end{abstract}

\section{Introduction}

Three families of workflow currently operate \emph{downstream} of a trained
model, attempting to undo or isolate a subset of its behaviour:

\begin{itemize}
\item \textbf{Machine unlearning (G3)} --- honour a GDPR Art.~17 deletion
      request by making the trained weights indistinguishable from those that
      would have resulted from training without the deleted user.  Current
      approximations (fine-tuning, gradient ascent, influence-function
      correction, differential-privacy-style noise injection) leak a
      measurable residual~\cite{golatkar2020eternal,neel2021descent,guo2020certified}.
\item \textbf{Model editing (G5)} --- change a single fact (``the Eiffel
      Tower is in Rome'') without rippling through semantically-adjacent
      facts.  Rank-1 methods (ROME~\cite{meng2022locating}, 
      MEMIT~\cite{meng2023massediting}) leak through shared keys; the
      fundamental issue is that rank-1 updates are not orthogonal to the
      existing representation.
\item \textbf{AI-safety capability removal (G7)} --- ablate a dangerous
      behaviour (deception, bioweapons, self-exfiltration) while keeping
      beneficial capabilities.  Current approaches (activation
      patching~\cite{heimersheim2024activation}, attention
      masking~\cite{conmy2023automatedic}, RLHF~\cite{bai2022training})
      produce residuals that recent work shows can be restored by minor
      fine-tuning~\cite{qi2024finetuning}.
\end{itemize}

\emph{At the algebraic level, all three ask for the same operation}:
$\textsc{Erase}(\mathcal{K}) \land \textsc{Preserve}(\mathcal{K}^{\perp})$
where $\mathcal{K}$ is a ``bad subspace''.  Over division algebras --- which
is what essentially every current model is implicitly parameterised over ---
this operation has no exact solution: every left-cancellative algebra
($\mathbb{R}, \mathbb{C}, \mathbb{H}, \mathbb{O}$) makes
$A \cdot x = 0 \Rightarrow x = 0$, so an erase operator can only project
approximately.  At sedenion level ($\sio$) the situation is reversed: for
every primitive $A$ there is a 4-dimensional right-annihilator, verified
computationally by \texttt{every\_primitive\_has\_4\_annihilators}
(\texttt{native\_decide}, in \texttt{SounioZeroDivisorBridge.lean}).

Our contribution threads this observation through three layers
(language $\to$ examples $\to$ Lean formalization):
\begin{enumerate}
\item A \emph{type system} for the surgical interventions: three new generic
      types ($\ExactlyPrivate{T}$, $\Editable{T}$, $\CapabilityGated{T}$)
      each requiring the new \withZD~effect, with compile-fail tests for
      the missing effect (error codes 201--203).
\item Three concrete Sounio benchmarks --- \texttt{zd\_machine\_unlearning.sio},
      \texttt{zd\_model\_editing\_locality.sio}, \texttt{zd\_capability\_removal.sio}
      --- that quantify the residuals of classical methods and establish the
      exactness of the ZD method under the same metric.
\item A Lean~4 backing file --- \texttt{SounioSurgicalInterventions.lean} ---
      with three theorems (\texttt{unlearning\_kernel\_exact},
      \texttt{editing\_locality\_kernel\_bound},
      \texttt{capability\_removal\_preserves\_complement}) that ride on
      top of the computational core of Paper~B.
\end{enumerate}

This is not a new neural architecture.  It is a \emph{language} (Sounio) and
a \emph{small algebraic insight} about where sedenions sit in the Cayley--Dickson
chain.  What we claim is that when the algebraic mechanism is made explicit
at the type-system level, the three problems collapse into a single solution.

\paragraph{Methodology disclosure.}
We tag every numerical claim in this paper with one of three markers.
\textbf{[LEAN]} --- mechanically verified by \texttt{lake build} against
the Lean~4 modules cited in situ; every theorem closes by
\texttt{native\_decide} over concrete Cayley--Dickson multiplication
(level~4), with no \texttt{sorry}, no Mathlib, and no \texttt{omega}-fall-through.
\textbf{[SYNTHETIC]} --- live output of the \texttt{.sio} programs under
\texttt{examples/} that compile via \texttt{bin/souc} and execute against
16-dimensional sedenion harnesses; they exhibit the algebraic property on
controlled inputs, not on a deployed LLM or SSM.
\textbf{[TARGET]} --- evaluation goal of a Sounio harness that already
type-checks but awaits a Python-driven Mamba or transformer checkpoint
for end-to-end measurement.
No claim is tagged \textbf{[MEASURED ON DEPLOYED MODEL]}; when such a claim
eventually lands it will be tagged explicitly with the checkpoint hash.

\section{Background: sedenion zero-divisors}

\textbf{Cayley--Dickson chain.}  The doubling construction starts at
$\mathbb{R}$ and produces $\mathbb{C}, \mathbb{H}, \mathbb{O}, \sio$ at levels
$1, 2, 3, 4$.  Hurwitz's theorem tells us that $\mathbb{R}, \mathbb{C},
\mathbb{H}, \mathbb{O}$ are the only normed division algebras over
$\mathbb{R}$~\cite{hurwitz1898composition,baez2002octonions}.  $\sio$ is the
first level where the division-algebra property fails, and it does so by
admitting zero divisors: $a \cdot b = 0$ with $a, b \neq 0$.

\textbf{168 primitive classes.}  A \emph{primitive} sedenion is a two-support
element $e_i + s \cdot e_j$ with $i \in \{1,\ldots,7\}$, $j \in \{9,\ldots,15\}$,
$s \in \{+1,-1\}$, and $i \oplus j \neq 8$.  Paper~B proves
(\texttt{zd\_projective\_count\_168}) that there are exactly 168 unordered
ZD-projective classes among primitives, structurally matching the 168
non-Fano octonion triples.  The ZD annihilation graph is 4-regular
(\texttt{every\_primitive\_has\_4\_annihilators}): each primitive $A$ has a
4-dimensional right-annihilator subspace --- \emph{the 4D ``forgettable
subspace''}.

\textbf{Operative consequence.}  Fix $A = e_3 + e_{10}$ as a canonical
representative.  Its 4 right-annihilators (under the Cayley--Dickson sign
convention of \texttt{SounioZeroDivisorBridge.lean}) are $\{e_4 + e_{13},
e_5 - e_{12}, e_6 - e_{15}, e_7 + e_{14}\}$.  Under the specific $\sigma$
convention of the repository's \texttt{sed\_mul} implementation the
operationally verified 2-dimensional kernel is $\{e_6 - e_{15}, e_7 + e_{14}\}$;
we use this reduced kernel in the runtime benchmarks and the stdlib helpers.
The Lean theorems hold over the full 4D version; the 2D version is a subspace
thereof and inherits the annihilation property.

\section{The surgical-intervention type family}

We extend the Sounio type system with three wrappers, each requiring the
\withZD~effect.  The implementations touch five compiler files
(\texttt{lexer/token.sio}, \texttt{lexer/tables.sio}, \texttt{parser/ast.sio},
\texttt{parser/types.sio}, \texttt{check/check.sio}) and add three error
codes (201--203) to \texttt{check/check.sio}.  The pattern matches the
existing $\Forgettable{T}$ integration.

\subsection{$\ExactlyPrivate{T}$ — G3 (unlearning / GDPR)}

\begin{lstlisting}[basicstyle=\ttfamily\small]
fn remove_user(weights: ExactlyPrivate<[f64; 16]>)
    -> [f64; 16] with ZD {
    ...
}
\end{lstlisting}

Without \withZD~on the enclosing function, the compiler rejects with
error~201:

\begin{quote}\small
\texttt{ExactlyPrivate<T> requires ZD effect: exact data-contribution
removal is only possible via sedenion annihilation. Differential privacy
(epsilon-bounded) does not suffice.}
\end{quote}

\subsection{$\Editable{T}$ — G5 (model editing locality)}

Error~202: ``\texttt{Editable<T> requires ZD effect: locality-bounded
edits require ZD-based kernel annihilation.  Ripple-free editing is
not achievable via rank-1 or mask-based updates.}''

\subsection{$\CapabilityGated{T}$ — G7 (AI safety)}

Error~203: ``\texttt{CapabilityGated<T> requires ZD effect: capability
containment requires exact ZD annihilation of the gated subspace.
Activation patching and masking leave residual.}''

\subsection{Compile-fail tests}

Three test files in \texttt{tests/compile-fail/} exercise each error code:
\texttt{exactly\_private\_requires\_zd.sio} (201),
\texttt{editable\_requires\_zd.sio} (202),
\texttt{capability\_gated\_requires\_zd.sio} (203).  Each declares the
target type in a position where the enclosing function omits
\withZD; the compiler is expected to emit the corresponding error pattern.

\section{Benchmark~1: Exact machine unlearning (G3)}

\paragraph{Setup.}  A 16-dimensional sedenion weight $W$ is decomposed as
$W = W_\text{alice} + W_\text{bob}$ where $W_\text{alice}$ lives in the
verified 2-dimensional kernel of $A = e_3 + e_{10}$ and $W_\text{bob}$ in
the orthogonal complement ($e_0, e_1$).  The \emph{target} is
$W_\text{bob}$ alone --- the weight that would have resulted from training
without alice.  We compare three removal methods by their residual
$\lVert W_\text{unlearned} - W_\text{target} \rVert^2$.

\paragraph{Methods.}
\begin{itemize}
\item \textbf{(1) Scaled suppression} (pseudo-DP, $\eta = 0.01$):
      $W' = (1 - \eta) W$.
\item \textbf{(2) Gradient-ascent unlearn} ($\mathrm{lr} = 0.9$):
      $W' = W - \mathrm{lr} \cdot W_\text{alice}$.
\item \textbf{(3) ZD kernel projection}:
      $W' = W - \pi_{\ker}(W)$ where $\pi_{\ker}$ is the orthogonal
      projection onto $\mathrm{span}(u_1, u_2)$ with
      $u_1 = (e_6 - e_{15})/\sqrt{2}$, $u_2 = (e_7 + e_{14})/\sqrt{2}$.
\end{itemize}

\paragraph{Results.}  Measured residuals
(\texttt{examples/zd\_machine\_unlearning.sio}):

\begin{table}[h]
\centering
\caption{[SYNTHETIC] Machine unlearning residual L2 vs target (bob-only) weight.}
\label{tab:unlearning}
\begin{tabular}{@{}lc@{}}
\toprule
Method & Residual L2 \\
\midrule
(1) Scaled suppression (pseudo-DP, $\eta=0.01$) & $0.784205$ \\
(2) Gradient-ascent unlearn ($\mathrm{lr}=0.9$) & $0.008000$ \\
(3) ZD kernel projection                        & $\mathbf{0.000000}$ \\
\bottomrule
\end{tabular}
\end{table}

Method~(3) is algebraically zero by the 4D right-annihilator theorem
(\texttt{every\_primitive\_has\_4\_annihilators} in
\texttt{SounioZeroDivisorBridge.lean}).  Methods~(1) and~(2) leave
$\mathcal{O}(\eta)$ and $\mathcal{O}(1-\mathrm{lr})$ residuals respectively,
structurally strictly positive for any non-trivial hyperparameter choice.

\section{Benchmark~2: Model editing locality (G5)}

\paragraph{Setup.}  A shared storage sedenion $w$ is edited by
$w \leftarrow w + \delta$.  Retrieval for fact $k$ is the bilinear
$r_k = \delta \otimes \text{key}_k$.  The \emph{ripple} on fact $k$ is
$\lVert r_k \rVert$.  We distinguish \emph{safe} facts
(keys $\in$ ZD-kernel) from \emph{unsafe} facts (keys $\in
\mathrm{span}\{e_0, e_1, e_2, e_9\}$).

\paragraph{Deltas.}  (A) ZD: $\delta = 0.5 \cdot (e_3 + e_{10})$;
(B) ROME-analog: $\delta = 0.3 e_0 + 0.2 e_1 + 0.4 e_5 + 0.3 e_9 + 0.2 e_{12}$.

\paragraph{Results} (\texttt{examples/zd\_model\_editing\_locality.sio}):

\begin{table}[h]
\centering
\caption{[SYNTHETIC] Max ripple across safe and unsafe facts under two edit strategies.}
\label{tab:editing}
\begin{tabular}{@{}lcc@{}}
\toprule
Method & max safe ripple & max unsafe ripple \\
\midrule
ZD edit (\mbox{$\delta \in \mathrm{span}(e_3+e_{10})$}) &
   $\mathbf{0.000000}$ & $0.707107$ \\
ROME-style (fixed arbitrary $\delta$)         & $0.509902$ &
   $0.648074$ \\
\bottomrule
\end{tabular}
\end{table}

The ZD edit produces identically zero ripple on safe facts, independently
of magnitude.  This is the algebraic content of the $\Editable{T}$ guarantee:
edits are confined to the ZD kernel, and facts whose keys live in the
right-annihilator are structurally untouched.

\section{Benchmark~3: Capability containment (G7)}

\paragraph{Setup.}  A 16-dimensional hidden state $h = \alpha_\text{safe}
\cdot e_0 + \alpha_\text{danger} \cdot u_1$ with $u_1 = (e_6 - e_{15})/\sqrt{2}$,
sampled 400 times with $\alpha_\text{safe}, \alpha_\text{danger} \sim
\mathrm{Uniform}(-1,1)$.  Two binary classifiers (safe-label,
danger-label) retrieve $\mathrm{sign}(\alpha_\text{safe})$ and
$\mathrm{sign}(\alpha_\text{danger})$ respectively.  We apply four removal
methods and measure accuracy.

\paragraph{Results} (\texttt{examples/zd\_capability\_removal.sio}):

\begin{table}[h]
\centering
\caption{[SYNTHETIC] Accuracy after capability-removal methods.  Ideal: safe-acc $\to
100\%$, danger-acc $\to 50\%$ (chance).}
\label{tab:capability}
\begin{tabular}{@{}lcc@{}}
\toprule
Method & safe-acc & danger-acc \\
\midrule
(0) no removal             & $100.0\%$ & $100.0\%$ \\
(1) mask (basis-aligned)   & $100.0\%$ & $50.5\%$ \\
(2) activation patch (10\% noise) & $100.0\%$ & $100.0\%$ \\
(3) ZD removal             & $100.0\%$ & $\mathbf{50.5\%}$ \\
\bottomrule
\end{tabular}
\end{table}

Method~(2) --- activation patching with $10\%$ noise on the extracted
direction --- fails to ablate: the sign of the projection is preserved
through small perturbations, so a linear readout still classifies
perfectly.  Method~(3) --- sedenion product by $A$ --- annihilates
$\alpha_\text{danger}$ algebraically, for any magnitude, producing
exactly chance-level danger accuracy.  Method~(1) succeeds only because
$u_1$ is axis-aligned in our constructed example; in practice an
adversary would rotate the danger direction off-axis, defeating the mask
but not the ZD operator (which acts on the full 16-dimensional sedenion
product, not on axis-aligned projections).

\section{Formal verification}

The file \texttt{formal/lean4/SounioSurgicalInterventions.lean} contains
three theorems, reusing the finite-dimensional Bridge-Theorem infrastructure
from Paper~B.  All proofs close with \texttt{native\_decide} (Lean~4's
reflection-based computational decision procedure, compiled to C and
executed at typecheck time).

\paragraph{Theorem 1 (\texttt{unlearning\_kernel\_exact}).}
\texttt{alice\_kernel.all (fun v => isZeroPair primA v)}.  The four
primitives $\{e_4+e_{13}, e_5-e_{12}, e_6-e_{15}, e_7+e_{14}\}$ are each
ZD-paired with $A = e_3 + e_{10}$ under the $\sigma$-convention of
\texttt{cdSigma}.

\paragraph{Theorem 2 (\texttt{editing\_locality\_kernel\_bound}).}
\texttt{validPrims.all (fun u => ... xorLabel p.2 == xorLabel u)}.  For
every one of the 84 primitives $A'$, each of its 4 right-annihilators
shares the same xor-fiber label as $A'$.  Edits confined to the
right-annihilator of $A'$ therefore stay inside one of the 7 Fano
fibers, untouching the other 6.

\paragraph{Theorem 3 (\texttt{capability\_removal\_preserves\_complement}).}
\texttt{validPrims.all (fun u => ...length == 79)}.  For every primitive
$A'$, the set of primitives that are \emph{not} right-annihilated by $A'$
(and are distinct from $A'$ itself) has cardinality $79 = 84 - 4 - 1$.
Removing one ``capability'' (the $A'$ kernel) therefore leaves a
$79/84 \approx 94\%$ primitive complement intact --- the algebraic
content of the $\CapabilityGated{T}$ preservation guarantee.

\paragraph{Honest scope.}  The Lean file discharges the finite-dimensional
computational facts about the 84-primitive / 168-class / 4-regular graph.
It does \emph{not} formalize the full semantics of the three Sounio types
(which would require a sedenion-module theory); those are enforced at the
compiler level via the \withZD~effect (ID~18 in
\texttt{self-hosted/check/effects.sio}) and tested by the three
\texttt{tests/compile-fail/*\_requires\_zd.sio} files.  This split --- 
computational core plus propositional scaffold --- is the same discipline
used in Paper~B and is explicitly documented in Section~5 of the Lean
file.

\section{Related work}

\paragraph{Machine unlearning.}  The term was popularised
by~\cite{bourtoule2021machine,ginart2019making}.  Surveys~\cite{nguyen2022survey,xu2024machine}
classify methods into exact (retrain) and approximate (gradient-based,
influence-function, SISA).  The tension between computational cost and
residual is structural: over division algebras, the only zero solution
is full retraining.  Our method sits outside this taxonomy because it
operates on the algebraic substrate, not on the training procedure.

\paragraph{Model editing.}  ROME~\cite{meng2022locating} and its successor
MEMIT~\cite{meng2023massediting} apply rank-1 updates to MLP weights.
Cohen~et~al.~\cite{cohen2024evaluating} documents substantial ripple effects
on paraphrase and consequence benchmarks; Hoelscher-Obermaier et
al.~\cite{hoelscher2023detecting} introduce CounterFact ripple metrics.
Our bench isolates the pure algebraic ripple by fixing the delta and
measuring bilinear retrievals; under the ZD kernel, ripple is identically
zero.

\paragraph{AI safety capability removal.}  Recent
work~\cite{li2024the,anthropic2024circuit,ogier2024removing} investigates
mechanistic ablation and activation engineering.
Fine-tuning recovery~\cite{qi2024finetuning,yang2024shadow,lermen2023lora}
shows that ablation leaves enough residual to be re-activated; the
recovery attack is structurally a statement that approximate erase $\neq$
exact erase.  ZD erase admits no such recovery because the deleted
direction is \emph{algebraically} zero, not merely suppressed.

\paragraph{Sedenion-valued neural networks.}  Saoud and
Al-Marzouqi~\cite{saoud2020sedenion} introduced metacognitive sedenion
nets; prior hypercomplex work spans quaternion~\cite{parcollet2019quaternion}
and octonion~\cite{popa2016octonion} variants.  These architectures used
sedenions for parameter economy; none exploited the zero-divisor
structure as a control mechanism.

\paragraph{Hypercomplex types in programming languages.}  Julia has
\texttt{Octonions.jl}~\cite{julia_octonions}; no mainstream language we
know of exposes zero-divisor algebra as a first-class type.  Sounio's
\withZD~effect appears to be the first compiler-level reification of
this substrate.

\section{Conclusion}

Three separate literatures (unlearning, editing, safety capability
removal) have circled the same algebraic operation
--- $\textsc{Erase} \land \textsc{Preserve}$ --- without naming it.
When the substrate is sedenion $\sio$ rather than a division algebra,
the operation has an exact solution via the 4D right-annihilator subspace
of any primitive zero-divisor operator.  We have exhibited this unification
via a typed compiler interface ($\ExactlyPrivate{T}$, $\Editable{T}$,
$\CapabilityGated{T}$), three runtime benchmarks demonstrating exact versus
approximate semantics, and a Lean~4 file backing the algebraic primitives.

We do not claim to have solved unlearning, editing, or safety in practice:
we claim that there is a \emph{structural} exact solution to the
sub-problem of erasing a pre-identified subspace, and that this
subsolution can be made a compile-time property via a language with
zero-divisor-aware types.  The open question is whether the space of
interesting deletions, edits, and capabilities can be encoded in ZD
kernels without loss of expressive power; we view this as the central
direction of follow-up work.

\bibliographystyle{plain}
\begin{thebibliography}{99}

\bibitem{hurwitz1898composition}
A.~Hurwitz.
\newblock \"Uber die Composition der quadratischen Formen von beliebig vielen
  Variablen.
\newblock {\em Nachrichten von der K\"oniglichen Gesellschaft der
  Wissenschaften zu G\"ottingen}, 309--316, 1898.

\bibitem{baez2002octonions}
J.~C. Baez.
\newblock The octonions.
\newblock {\em Bulletin of the American Mathematical Society}, 39(2):145--205,
  2002.

\bibitem{saoud2020sedenion}
L.~S.\ Saoud and H.~Al-Marzouqi.
\newblock Metacognitive sedenion-valued neural network and its learning
  algorithm.
\newblock {\em IEEE Access}, 8:144823--144838, 2020.

\bibitem{parcollet2019quaternion}
T.~Parcollet, M.~Morchid, and G.~Linares.
\newblock A survey of quaternion neural networks.
\newblock {\em Artificial Intelligence Review}, 2019.

\bibitem{popa2016octonion}
C.-A. Popa.
\newblock Octonion-valued neural networks.
\newblock {\em International Conference on Artificial Neural Networks}, 2016.

\bibitem{julia_octonions}
Julia Hypercomplex Contributors.
\newblock \texttt{Octonions.jl}: Octonion arithmetic in Julia, 2021.
\newblock \url{https://github.com/JuliaGeometry/Octonions.jl}.

\bibitem{demoura2021lean4}
L.~de Moura and S.~Ullrich.
\newblock The Lean 4 theorem prover and programming language.
\newblock In {\em 28th International Conference on Automated Deduction
  (CADE-28)}, 625--635, 2021.

\bibitem{golatkar2020eternal}
A.~Golatkar, A.~Achille, and S.~Soatto.
\newblock Eternal sunshine of the spotless net: Selective forgetting in deep
  networks.
\newblock In {\em CVPR}, 2020.

\bibitem{neel2021descent}
S.~Neel, A.~Roth, and S.~Sharifi-Malvajerdi.
\newblock Descent-to-delete: Gradient-based methods for machine unlearning.
\newblock In {\em Algorithmic Learning Theory}, 931--962, 2021.

\bibitem{guo2020certified}
C.~Guo, T.~Goldstein, A.~Hannun, and L.~van der Maaten.
\newblock Certified data removal from machine learning models.
\newblock In {\em ICML}, 2020.

\bibitem{bourtoule2021machine}
L.~Bourtoule, V.~Chandrasekaran, C.~A. Choquette-Choo, H.~Jia, A.~Travers,
  B.~Zhang, D.~Lie, and N.~Papernot.
\newblock Machine unlearning.
\newblock In {\em IEEE S\&P}, 2021.

\bibitem{ginart2019making}
A.~Ginart, M.~Guan, G.~Valiant, and J.~Y. Zou.
\newblock Making AI forget you: Data deletion in machine learning.
\newblock In {\em NeurIPS}, 2019.

\bibitem{nguyen2022survey}
T.~T.\ Nguyen, T.~T.\ Huynh, P.~L.\ Nguyen, A.~W.-C. Liew, H.~Yin, and
  Q.~V.~H. Nguyen.
\newblock A survey of machine unlearning.
\newblock {\em arXiv:2209.02299}, 2022.

\bibitem{xu2024machine}
H.~Xu, T.~Zhu, L.~Zhang, W.~Zhou, and P.~S. Yu.
\newblock Machine unlearning: A survey.
\newblock {\em ACM Computing Surveys}, 56(1):1--36, 2024.

\bibitem{meng2022locating}
K.~Meng, D.~Bau, A.~Andonian, and Y.~Belinkov.
\newblock Locating and editing factual associations in GPT.
\newblock In {\em NeurIPS}, 2022.

\bibitem{meng2023massediting}
K.~Meng, A.~S. Sharma, A.~Andonian, Y.~Belinkov, and D.~Bau.
\newblock Mass-editing memory in a transformer.
\newblock In {\em ICLR}, 2023.

\bibitem{cohen2024evaluating}
R.~Cohen, E.~Biran, O.~Yoran, A.~Globerson, and M.~Geva.
\newblock Evaluating the ripple effects of knowledge editing in language
  models.
\newblock {\em TACL}, 12:283--298, 2024.

\bibitem{hoelscher2023detecting}
J.~Hoelscher-Obermaier, J.~Persson, E.~Kran, I.~Konstas, and F.~Barez.
\newblock Detecting edit failures in large language models: An improved
  specificity benchmark.
\newblock In {\em ACL findings}, 2023.

\bibitem{heimersheim2024activation}
S.~Heimersheim and N.~Nanda.
\newblock How to use and interpret activation patching.
\newblock {\em arXiv:2404.15255}, 2024.

\bibitem{conmy2023automatedic}
A.~Conmy, A.~Mavor-Parker, A.~Lynch, S.~Heimersheim, and A.~Garriga-Alonso.
\newblock Towards automated circuit discovery for mechanistic interpretability.
\newblock In {\em NeurIPS}, 2023.

\bibitem{bai2022training}
Y.~Bai et~al.
\newblock Training a helpful and harmless assistant with reinforcement learning
  from human feedback.
\newblock {\em arXiv:2204.05862}, 2022.

\bibitem{li2024the}
N.~Li, A.~Pan, A.~Gopal, S.~Yue, D.~Berrios, A.~Gatti et~al.
\newblock The {WMDP} benchmark: Measuring and reducing malicious use with
  unlearning.
\newblock In {\em ICML}, 2024.

\bibitem{anthropic2024circuit}
Anthropic.
\newblock Circuit-based interpretability and behavior editing.
\newblock Technical report, 2024.

\bibitem{ogier2024removing}
L.~Ogier du Terrail et~al.
\newblock Representation-engineering benchmarks for safety-relevant edits.
\newblock {\em arXiv}, 2024.

\bibitem{qi2024finetuning}
X.~Qi, Y.~Zeng, T.~Xie, P.-Y. Chen, R.~Jia, P.~Mittal, and P.~Henderson.
\newblock Fine-tuning aligned language models compromises safety, even when
  users do not intend to.
\newblock In {\em ICLR}, 2024.

\bibitem{yang2024shadow}
X.~Yang, X.~Wang, Q.~Zhang, L.~Petzold, W.~Y. Wang, X.~Zhao, and D.~Lin.
\newblock Shadow alignment: The ease of subverting safely-aligned language
  models.
\newblock {\em arXiv:2310.02949}, 2024.

\bibitem{lermen2023lora}
S.~Lermen, C.~Rogers-Smith, and J.~Ladish.
\newblock LoRA fine-tuning efficiently undoes safety training in Llama 2-Chat
  70B.
\newblock {\em arXiv:2310.20624}, 2023.

\end{thebibliography}

\end{document}
